\section{IGSSE}
\label{sec:method}

\subsection{Problem boundary and representation target}
\label{sec:setup}
Let $G=(V,E,\mathbf X)$ be an attributed graph with adjacency $\mathbf A$ and
$n$ nodes.  The task label $y$ is a node class, graph label, or continuous
target, but never enters the embedding path.  \method{} produces
$\Hset\in\mathbb R^{n\times d}$, where row $i$ remains associated with node
$i$.  This row correspondence is essential: a content decoder must read a node
target from one point, while a relation decoder must read a graph relation from a
pair of points.

We distinguish the embedding target from the prediction target.  The recovery
target is $T(G)=(\mathbf A,\mathbf C)$, where
$\mathbf C=\operatorname{sg}(\operatorname{Enc}_{X}(\mathbf X))$ and
$\operatorname{sg}$ stops gradients.  A configuration may instead recover a
standardized raw feature target.  In either case, labels are excluded.  The
stopped encoded target prevents the feature decoder from reducing reconstruction
to a shortcut through the target encoder.  We reserve $\Hset_{\mathrm{dec}}$
for the final state audited by the two decoders, and reserve
$\Hset_{\mathrm{task}}$ for a final-state or jump-concatenated reader input.
This separation keeps a task reader from silently changing the object constrained
by recovery.

\paragraph{Core-path correspondence.}
The core computation first forms original coordinates $\mathbf h_{o,i}$; source
queries, target keys, and source values produce $e_{si}$ and $\alpha_{si}$;
their weighted value sum forms the offset $\mathbf o_i$; and the offset is fused
with the original coordinate before the feed-forward map produces $\mathbf h_i$.
The first-layer path is
\begin{equation}
 \mathbf x_i\rightarrow\mathbf h_{o,i}\rightarrow
 (\mathbf q_s,\mathbf k_i,\mathbf v_s)\rightarrow e_{si}\rightarrow
 \alpha_{si}\rightarrow\mathbf o_i\rightarrow
 \mathbf r_i=\mathbf h_{o,i}+\mathbf o_i\rightarrow\mathbf h_i.
 \label{eq:visual-core}
\end{equation}
Later layers repeat this path with the preceding point state in place of
$\mathbf h_o$.  Noise is a training configuration: it perturbs attention
energies but does not alter this deterministic inference path.

\subsection{Feature--structure deviation encoding}
\label{sec:fsde}
First, a shared attribute map forms an edge-free baseline:
\begin{equation}
  \mathbf h^{(0)}_i=\operatorname{Enc}_X(\mathbf x_i).
  \label{eq:baseline}
\end{equation}
In the default implementation, $\operatorname{Enc}_X$ applies input layer
normalization, a linear map to a hidden space, GELU and dropout, a second linear
map, and output normalization.  Formally,
\begin{align}
 \widetilde{\mathbf x}_i&=\operatorname{LN}_{\mathrm{in}}(\mathbf x_i),\\
 \mathbf z_i&=\operatorname{Dropout}\bigl(
 \operatorname{GELU}(\widetilde{\mathbf x}_i\mathbf W_1+\mathbf b_1)\bigr),\\
 \mathbf h^{(0)}_i&=\operatorname{LN}_{\mathrm{out}}(\mathbf z_i\mathbf W_2+\mathbf b_2).
 \label{eq:feature-encoder}
\end{align}
This encoder has no access to $E$, so its output provides an explicit content
reference for the subsequent structural correction and for the content target.
The baseline gives a concrete location for node content before the graph changes
it.  At layer $\ell$, the encoding edge set includes self loops,
$E_{\mathrm{enc}}=E\cup\{(i,i):i\in V\}$.  Multi-head projections produce
$\mathbf q_i$, $\mathbf k_i$, and $\mathbf v_i$.  For source $s$, target $i$,
and head $m$, \method{} computes
\begin{equation}
 e_{si}^{(m)}=\operatorname{LeakyReLU}\left(
 (\mathbf a_{\mathrm{src}}^{(m)})^\top\mathbf q_s^{(m)}+
 (\mathbf a_{\mathrm{dst}}^{(m)})^\top\mathbf k_i^{(m)}\right).
 \label{eq:energy}
\end{equation}

\paragraph{Noise curriculum.}
The noise network $\psi$ receives the mean head energy and a normalized training
progress $u$.  It predicts a nonnegative scale $\sigma_{si}^{(m)}$, and training
uses
\begin{equation}
 \widetilde e_{si}^{(m)}=e_{si}^{(m)}+
 \sigma_{si}^{(m)}\max(0,1-u)\epsilon_{si}^{(m)},\qquad
 \epsilon_{si}^{(m)}\sim\mathcal N(0,1).
 \label{eq:noise}
\end{equation}
More specifically, the mean energy $\bar e_{si}=M^{-1}\sum_m e_{si}^{(m)}$
and progress $u$ enter a two-layer MLP.  Its output is added to a learned global
log-scale, transformed by softplus, lower bounded for numerical stability, and
clipped to a configured maximum.  Each head therefore receives its own
edge-dependent perturbation scale.  The survival term in Eq.~\eqref{eq:objective}
prevents this scale from trivially vanishing during the noise phase.
Evaluation sets $\widetilde e=e$.  The curriculum perturbs early relation
weights, then removes the perturbation before selection and test evaluation.
The noise is a training configuration rather than a separate graph encoder.

For each target, a softmax over incoming encoding edges gives
$\alpha_{si}^{(m)}$.  The structural deviation is the projected multi-head sum
\begin{equation}
 \mathbf o_i=\mathbf W_O\left\Vert_m
 \sum_{s:(s,i)\in E_{\mathrm{enc}}}\alpha_{si}^{(m)}\mathbf v_s^{(m)}
 \right\Vert.
 \label{eq:deviation}
\end{equation}
Residual fusion then updates the point state:
\begin{align}
 \mathbf r_i^{(\ell)}&=\mathbf h_i^{(\ell-1)}+
 \operatorname{Dropout}(\mathbf o_i^{(\ell)}),\\
 \widehat{\mathbf h}_i^{(\ell)}&=\operatorname{LN}(\mathbf r_i^{(\ell)}),\\
 \mathbf h_i^{(\ell)}&=\widehat{\mathbf h}_i^{(\ell)}+
 \operatorname{FFN}(\widehat{\mathbf h}_i^{(\ell)}).
 \label{eq:fusion}
\end{align}
After $L$ layers, $\Hset_{\mathrm{dec}}=\mathbf H^{(L)}$ is used for recovery.
Node tasks may use this final state or a jump-concatenated state
$\mathbf h^{(0)}_i\Vert\cdots\Vert\mathbf h^{(L)}_i$.

\paragraph{Pseudocode 1: FSDE encoding.}
\begin{quote}\small
\textbf{Input:} $\mathbf X$, $E_{\mathrm{enc}}$, and progress $u$.\\
\textbf{Output:} $\Hset_{\mathrm{dec}}$, $\Hset_{\mathrm{task}}$, and
mean noise scale $\overline\sigma$.\\
\textbf{1.} Set $\mathbf H^{(0)}\leftarrow\operatorname{Enc}_X(\mathbf X)$.\\
\textbf{2.} For each layer, project points to queries, keys, and values. For
each source--target edge $(s,i)$, compute $e_{si}$, normalize incoming scores,
and sum $\alpha_{si}\mathbf v_s$ into $\mathbf o_i$.\\
\textbf{3.} Perturb energies with Eq.~\eqref{eq:noise} during training; use
clean energies otherwise. Form $\mathbf r_i=\mathbf h_i^{(\ell-1)}+
\mathbf o_i$ and apply the normalized feed-forward fusion in Eq.~\eqref{eq:fusion}.\\
\textbf{4.} Return $\Hset_{\mathrm{dec}}\leftarrow\mathbf H^{(L)}$ and the
configured final or jump-concatenated task state.
\end{quote}

\subsection{Two recovery paths}
\label{sec:recovery}
The pair decoder scores sampled positive and negative pairs:
\begin{equation}
 \begin{aligned}
 \widehat A_{ij}&=\sigma(D_A(\mathbf h_i,\mathbf h_j)),\\
 \mathcal L_{\mathrm{link}}&=
 \operatorname{BCEWithLogits}(D_A(\mathbf h_i,\mathbf h_j),A_{ij}).
 \end{aligned}
 \label{eq:link}
\end{equation}
The implementation offers an MLP score on concatenated point features,
differences, and products, or an attention-based pair score.  Both use the same
link supervision interface.  A separate content decoder gives
\begin{equation}
 \widehat{\mathbf C}=D_C(\Hset_{\mathrm{dec}}),\qquad
 \mathcal L_{\mathrm{feat}}=\operatorname{MSE}(\widehat{\mathbf C},\mathbf C).
 \label{eq:feat}
\end{equation}
The two losses impose different requirements: link recovery tests the relation
geometry of point pairs, whereas feature recovery tests the information readable
from an individual point.  We call their joint evaluation the \emph{recovery
gate}.  A task score cannot replace this gate because a head may exploit a small
subset of a representation.

\paragraph{Pair sampling and target hygiene.}
Materializing every non-edge is unnecessary and can dominate the objective on a
sparse graph.  At each update, we sample positive pairs from the complete input
edge set and sample an equal-size set of non-edges, excluding self loops and
known positives.  If edge dropout is active in the encoder, it affects only the
computation graph $E_{\mathrm{enc}}$; link positives still come from the
un-dropped input graph.  This distinguishes an input perturbation from a changed
reconstruction target.  The decoder therefore learns whether a point set
constructed under perturbation can still distinguish the declared relations.

The content target has the same hygiene requirement.  In the default
configuration, $\mathbf C=\operatorname{sg}(\operatorname{Enc}_X(\mathbf X))$
uses clean attributes even when the encoder receives an input perturbation.  A
direct-feature configuration replaces $\mathbf C$ by standardized $\mathbf X$.
Both variants preserve the label boundary: $y$ appears only in
$\mathcal L_{\mathrm{pred}}$ over the training mask.  This boundary prevents
the recovery path from observing the target that a downstream task is meant to
predict.

\paragraph{Decoder parameterizations.}
The default pair decoder is an MLP over complementary point-pair features,
\begin{equation}
 s_{ij}=\operatorname{MLP}_{\mathrm{link}}
 \left([\mathbf h_i;\mathbf h_j;\mathbf h_i-\mathbf h_j;
 \mathbf h_i\odot\mathbf h_j]\right).
 \label{eq:pairmlp}
\end{equation}
Concatenation preserves endpoint-specific content, the difference exposes a
directed displacement, and the product exposes coordinate agreement.  The MLP
returns a logit, so the sampling distribution rather than a fixed threshold
defines the training objective.  For undirected graphs, evaluation scores both
orientations or uses the configured symmetrization rule.

An attention-based decoder provides the alternate parameterization.  It projects
each point into head-specific queries and keys, scores their scaled inner product,
averages the two endpoint directions when the relation is undirected, and applies
a learned scale and bias.  The choice changes the relation function but not the
declared target, pair sampler, loss, or recovery metrics.  We therefore treat
decoder choice as a relation-geometry ablation rather than as a different
embedding problem.

For completeness, the attention score for an undirected candidate can be written
as
\begin{equation}
 \begin{aligned}
 s_{ij}&=\left(\operatorname{softplus}(\gamma)+0.5\right)
 \frac{1}{2M}\sum_m\left[
 \frac{(\mathbf q_i^{(m)})^\top\mathbf k_j^{(m)}}{\sqrt{d_h}}\right.\\
 &\hspace{18mm}\left.+
 \frac{(\mathbf q_j^{(m)})^\top\mathbf k_i^{(m)}}{\sqrt{d_h}}
 \right]+\beta .
 \end{aligned}
 \label{eq:pairattn}
\end{equation}
Here $\gamma$ and $\beta$ calibrate score scale and offset.  This decoder is
permutation equivariant because point rows, queries, and keys permute together.
It does not use a fixed adjacency after $\Hset_{\mathrm{dec}}$ is constructed.

The content decoder is a layer-normalized pointwise MLP.  It cannot aggregate
neighbors after $\Hset_{\mathrm{dec}}$ has been formed.  Feature reconstruction
therefore tests information readable from each final point, rather than a second
message-passing computation over the graph.  When edge attributes are required,
the interface extends with $D_E(\mathbf h_i,\mathbf h_j)$ and an explicit
edge-attribute loss; no current result claims that unmodeled edge attributes are
preserved.

\subsection{Training schedule and evaluation protocol}
\label{sec:training}
The optimization schedule separates three quantities that otherwise become easy
to conflate: a noisy encoder state during training, a clean deterministic state
for validation and test, and a task label visible only on training examples.
For a node task, the normalized progress is
\begin{equation}
 u=\min\left(1,\frac{\mathrm{epoch}}{0.4\,\mathrm{epochs}}\right).
 \label{eq:curriculum}
\end{equation}
The first 40\% of training therefore anneals Eq.~\eqref{eq:noise}; later
updates and all selection passes use unperturbed relation energies.  Gradient
clipping bounds the joint update, and the validation metric chooses a checkpoint.
The test split is read once from the selected checkpoint.  The default node-task
configuration uses dimension 64, three encoder layers, four heads, hidden width
256, dropout 0.4, Adam at $10^{-3}$ with weight decay $5\times10^{-4}$, loss
weights $(1.0,0.5,0.3,0.02)$, and up to 4096 positive and 4096 negative pairs.

\paragraph{Pseudocode 2: one joint update.}
\begin{quote}\small
\textbf{1.} Compute $u$, optionally perturb $\mathbf X$ and the encoder edge
view, then run Pseudocode~1 with training noise enabled.\\
\textbf{2.} Decode clean content targets from $\Hset_{\mathrm{dec}}$ and compute
$\mathcal L_{\mathrm{feat}}$. Sample pairs from the complete input graph and
compute $\mathcal L_{\mathrm{link}}$.\\
\textbf{3.} Read $\Hset_{\mathrm{task}}$ with the configured node or graph head;
compute $\mathcal L_{\mathrm{pred}}$ only on training labels.\\
\textbf{4.} Add $\mathcal L_{\mathrm{alive}}$, backpropagate the weighted sum,
clip gradients, and update encoder, decoders, and reader jointly.\\
\textbf{5.} For selection and test, disable noise, set $u=1$, and evaluate the
selected checkpoint.
\end{quote}

The reported embedding probe and direct recovery metrics have different roles.
The probe asks whether a standardized structural feature can predict links from
the learned state.  The direct recovery gate instead evaluates $D_A$ and $D_C$.
The paper reports them separately: a high probe score cannot certify feature
recovery, and a high task score cannot certify either target.  This distinction
also determines the experimental order: a data family first passes the recovery
gate, then opens a point or set prediction comparison.

\subsection{Loss attachment and gradient routes}
\label{sec:losses}
The joint objective in Eq.~\eqref{eq:objective} has four distinct attachment
points.  The prediction term supervises a task reader.  The link term sends
gradients through the pair decoder and the final point set.  The feature term
sends gradients through the content decoder and final point set, but not through
the stopped target encoder.  The survival term reaches the scale network and the
attention-energy path.  Thus, the representation is not trained by a single
``reconstruction loss'': relation geometry, node content, task relevance, and
the training perturbation receive separate signals.  Only
$\mathcal L_{\mathrm{pred}}$ reads $y$ on the training mask; the remaining
three terms use no task label.

For an embedding-only run, $\lambda_{pred}=0$ and the three representation terms
train a point set without a task label.  A task run restores
$\mathcal L_{pred}$ but retains the other terms.  This schedule makes two
experimental modes explicit: a recovery gate assesses the representation alone,
and a joint task run assesses the representation-reader system while keeping the
declared information constraints active.

\subsection{Point, set, and growth interfaces}
\label{sec:interfaces}
A shared point head maps each $\mathbf h_i$ to a node label or value.  A graph
task uses invariant attention pooling,
\begin{equation}
 \mathbf g=\sum_i\beta_i\mathbf U_V\mathbf h_i,\qquad
 \widehat Y_G=R_{\mathrm{graph}}(\mathbf g),
 \label{eq:pool}
\end{equation}
where $\beta_i$ is normalized over the graph.  The readout receives points only;
it does not reopen the original adjacency as a task message-passing graph.  The
frozen-$\Hset$ and readout-replacement ablations isolate this interface choice.

\paragraph{Output-level choice.}
For node classification and regression, the head is shared over point rows and
preserves node alignment.  When a node needs task context, a set-attention block
updates the entire point set before the shared head reads its row.  For graph
classification and regression, attention pooling or PMA maps a variable-cardinality
set to one invariant state.  Mean pooling is included only as an ablation because
it removes learned selection among points.  These output choices do not alter the
recovery target, so a change in task score can be attributed to the reader of
$\Hset$ rather than a different embedding object.

\paragraph{Concrete readers.}
For node classification, a shared MLP returns class logits
$\widehat{\mathbf z}_i=R_{\mathrm{node}}(\mathbf h_i)$ and optimizes masked
cross entropy.  Class weights and label smoothing, when configured, apply only
to training nodes.  For node regression, the same row-wise reader returns a
scalar or vector and uses Huber or squared error; validation selects the stated
regression metric.  These readers retain row alignment, so an input permutation
permutes their outputs in the same way.

For graph tasks, a learned query or attention weight reads a set state.  A PMA
configuration uses one or more seed queries; a simpler attention configuration
uses $\beta_i\propto\exp(\mathbf w^\top\mathbf h_i)$.  Both preserve invariance
to node ordering.  The graph head then maps $\mathbf g$ to a class distribution
or continuous property.  Graph-level recovery losses are averaged over graphs in
the batch, so a high graph score cannot silence the representation objective for
individual graphs.

For constrained autoregression, the reference construction first expands the
point set from $\Hset_s=\{\mathbf h_1,\ldots,\mathbf h_n\}$ to
$\Hset' =\{\mathbf h_1,\ldots,\mathbf h_n,\mathbf h_{n+1}\}$, then applies the
same pair-attention relation decoder to recover the expanded graph.  Formally,
\begin{equation}
 \Hset_{s+1}=\operatorname{StepSet}(\Hset_s,\boldsymbol\epsilon_s),
 \qquad E_{s+1}=\{(i,j):\sigma(D_A(\mathbf h_i,\mathbf h_j))\geq\tau\}.
 \label{eq:grow}
\end{equation}
Training may use a graph-family legality signal.  At evaluation, parameters are
fixed and the model free-grows without a step-wise rule teacher.  We report
terminal validity together with an embedding probe, because a legal output alone
does not establish that the state remains relationally readable.

\paragraph{State expansion, candidates, and live states.}
At a growth depth $s$, the expansion operator proposes one or more candidate
states from each live $\Hset_s$.  The relation decoder maps every candidate to
an edge score representation, and the graph-family predicate labels it legal or
illegal during training.  Candidate states can be ranked either by the declared
legality route (Legal) or by a learned continuation score (Net).  A live-state
capacity prevents the candidate tree from growing without bound; the capacity is
recorded as $\texttt{max\_live}$ in the generation protocol.

The distinction between construction and evaluation is deliberate.  Training
uses legality supervision to make the set transition sensitive to the graph
family.  Free-growth evaluation fixes all learned parameters, disables a
step-wise rule teacher, and evaluates terminal graphs at a larger target size.
Thus, the final legality predicate measures an output property while Emb F1
measures whether the resulting point state retains a usable relation geometry.
Neither quantity alone establishes the other.  The reported Legal and Net rows
therefore use the same target size and decoder interface but make their live-state
selection rule explicit.

\paragraph{Growth objective.}
For a candidate graph $G_s$ produced from $\Hset_s$, let
$c_s=\mathbb{I}[\operatorname{legal}(G_s)]$ be the family predicate.  A
legality head can optimize a binary loss $\mathcal L_{\mathrm{legal}}$ during
training, while the relation decoder continues to score the candidate graph.
The full generator objective combines the set-transition signal, pairwise
recovery, and legality signal according to the active graph-family protocol.
At test time, $c_s$ is never used as a step-wise teacher.  It is evaluated only
after the free-growth rollout to compute Valid.  This train--test separation is
what makes Legal and Net results evidence about a learned growth process rather
than a rule-filtered trajectory.

\paragraph{Pseudocode 3: constrained free growth.}
\begin{quote}\small
\textbf{Input:} learned expansion operator, pair decoder, initial set
$\Hset_0$, target size, gate type, and $\texttt{max\_live}$.\\
\textbf{1.} Initialize the live state collection with $\Hset_0$.\\
\textbf{2.} Until the target size, expand every live state with
$\operatorname{StepSet}$ and decode candidate relation scores.\\
\textbf{3.} During training, attach the graph-family legality signal to each
candidate. During evaluation, do not query the legality teacher.\\
\textbf{4.} Rank candidates by the declared Legal or Net gate, retain at most
$\texttt{max\_live}$ states, and continue with their point sets.\\
\textbf{5.} Decode terminal graphs and report Valid together with Emb F1 over
the completed rollouts.
\end{quote}

\subsection{Symmetry and objective}
\label{sec:objective}
Shared point maps and relation scoring preserve the intended permutation
structure.  For a consistent node permutation $\mathbf P$,
\begin{equation}
 \begin{aligned}
 \operatorname{Enc}(\mathbf P\mathbf A\mathbf P^\top,\mathbf P\mathbf X)
 &=\mathbf P\operatorname{Enc}(\mathbf A,\mathbf X),\\
 D_A(\mathbf P\Hset)&=\mathbf P D_A(\Hset)\mathbf P^\top,\\
 \operatorname{Pool}(\mathbf P\Hset)&=\operatorname{Pool}(\Hset).
 \end{aligned}
 \label{eq:symmetry}
\end{equation}
\paragraph{Why the symmetry holds.}
The attribute map is shared across rows, so it commutes with a node permutation.
For a fixed target, permuting the incoming edge order permutes the softmax
arguments and values but preserves their weighted sum.  Residual maps, layer
normalization, and pointwise feed-forward layers then preserve row alignment by
induction over encoder layers.  The pair decoder receives both point rows and
therefore transforms on both adjacency axes; attention pooling permutes weights
and values together, leaving the sum unchanged.  The argument assumes that any
edge dropout and noise sample are permuted consistently with the graph.

\paragraph{Recovery implication.}
Let $D$ be a deterministic decoder and let $d_T$ be a distance on declared
targets.  If $d_T(D(\operatorname{Enc}(G)),T(G))\leq\epsilon$ on the evaluation
domain, then two graphs $G_1,G_2$ with the same representation satisfy
\begin{equation}
 d_T(T(G_1),T(G_2))\leq 2\epsilon.
 \label{eq:recoverability}
\end{equation}
Indeed, both targets lie within $\epsilon$ of the same decoded representation,
so the result follows from the triangle inequality.  The statement explains why
direct recovery metrics matter: they bound how much declared input information a
shared embedding can merge.  It does not claim exact inversion at finite error,
reconstruction of unmodeled attributes, or a graph-isomorphism guarantee.
If a fixed decoder reconstructs the declared target with error at most
$\epsilon$, two graphs mapped to the same representation have target distance at
most $2\epsilon$ by the triangle inequality.  This is a recovery statement, not
a Weisfeiler--Lehman expressiveness claim.

Each task instance jointly optimizes
\begin{equation}
 \mathcal L=\lambda_{\mathrm{pred}}\mathcal L_{\mathrm{pred}}+
 \lambda_{\mathrm{link}}\mathcal L_{\mathrm{link}}+
 \lambda_{\mathrm{feat}}\mathcal L_{\mathrm{feat}}+
 \lambda_{\mathrm{alive}}\mathcal L_{\mathrm{alive}},
 \label{eq:objective}
\end{equation}
where $\mathcal L_{\mathrm{alive}}=-\log(\max(\overline\sigma,10^{-4}))$
keeps the learned training-noise scale active.  The prediction loss only accesses
training labels.  The recovery terms remain active in task mode, preserving an
embedding floor while the readout learns the task.

\paragraph{Computational boundary.}
The encoder aggregates only over $E_{\mathrm{enc}}$, so its layer cost is linear
in the number of encoded edges up to the head and hidden dimensions.  Relation
training uses sampled pairs, whereas a complete recovery evaluation may score a
larger specified pair set.  Set attention and unrestricted pair decoding can
become quadratic in the number of live points.  We therefore treat sampled-pair
training, chunked decoding, and live-state capacity as explicit protocol choices
rather than claiming graph-scale efficiency from small-graph results.  The
experiments report task accuracy and validity; measured time and memory remain a
separate evidence requirement.
